\documentclass[11pt]{article}
\newcommand{\ListaTitulo}{Gabarito --- Lista 04}
% Preâmbulo compartilhado — MATD48, Listas de Exercícios 2026
% Usado por Lista{NN}.tex e Gabarito{NN}.tex via % Preâmbulo compartilhado — MATD48, Listas de Exercícios 2026
% Usado por Lista{NN}.tex e Gabarito{NN}.tex via % Preâmbulo compartilhado — MATD48, Listas de Exercícios 2026
% Usado por Lista{NN}.tex e Gabarito{NN}.tex via \input{preamble}

\usepackage[utf8]{inputenc}
\usepackage[T1]{fontenc}
\usepackage[brazil]{babel}
\usepackage{fullpage}
\usepackage{amsmath,amssymb}
\usepackage{enumitem}
\usepackage{xcolor}
\usepackage{fancyhdr}
\usepackage{titlesec}
\usepackage{listings}
\usepackage{hyperref}
\usepackage{booktabs}
\usepackage{graphicx}

\definecolor{matd48azul}{HTML}{0B4F6C}
\definecolor{matd48verde}{HTML}{2E8B57}

\titleformat{\section}{\normalfont\Large\bfseries\color{matd48azul}}{\thesection}{1em}{}
\titleformat{\subsection}{\normalfont\large\bfseries\color{matd48azul}}{\thesubsection}{1em}{}

\providecommand{\ListaTitulo}{Lista de Exercícios}

\setlength{\headheight}{14pt}
\pagestyle{fancy}
\fancyhf{}
\lhead{\small MATD48 --- Planejamento de Experimentos A}
\rhead{\small \ListaTitulo}
\cfoot{\thepage}
\renewcommand{\headrulewidth}{0.4pt}

\lstset{
  basicstyle=\ttfamily\small,
  breaklines=true,
  frame=single,
  columns=fullflexible,
  backgroundcolor=\color{gray!8},
  commentstyle=\color{matd48verde},
  keywordstyle=\color{matd48azul}\bfseries,
  language=R,
  extendedchars=true,
  literate=%
    {á}{{\'a}}1 {à}{{\`a}}1 {â}{{\^a}}1 {ã}{{\~a}}1
    {é}{{\'e}}1 {ê}{{\^e}}1 {í}{{\'i}}1 {ó}{{\'o}}1
    {ô}{{\^o}}1 {õ}{{\~o}}1 {ú}{{\'u}}1 {ç}{{\c c}}1
    {Á}{{\'A}}1 {À}{{\`A}}1 {Â}{{\^A}}1 {Ã}{{\~A}}1
    {É}{{\'E}}1 {Ê}{{\^E}}1 {Í}{{\'I}}1 {Ó}{{\'O}}1
    {Ô}{{\^O}}1 {Õ}{{\~O}}1 {Ú}{{\'U}}1 {Ç}{{\c C}}1
}

\hypersetup{colorlinks=true, linkcolor=matd48azul, urlcolor=matd48azul, citecolor=matd48azul}

\newcommand{\cabecalho}[2]{%
  \begin{center}
    {\Large\bfseries\color{matd48azul} #1}\\[0.3em]
    {\large #2}\\[0.3em]
    \rule{\linewidth}{0.6pt}
  \end{center}
  \vspace{0.5em}
}

\newenvironment{questao}[1]{\vspace{1em}\noindent\textbf{Questão #1.}\hspace{0.4em}}{\par}
\newenvironment{solucao}{\vspace{0.3em}\noindent\textit{Solução.}\hspace{0.4em}\color{black!85}}{\par\vspace{0.5em}}


\usepackage[utf8]{inputenc}
\usepackage[T1]{fontenc}
\usepackage[brazil]{babel}
\usepackage{fullpage}
\usepackage{amsmath,amssymb}
\usepackage{enumitem}
\usepackage{xcolor}
\usepackage{fancyhdr}
\usepackage{titlesec}
\usepackage{listings}
\usepackage{hyperref}
\usepackage{booktabs}
\usepackage{graphicx}

\definecolor{matd48azul}{HTML}{0B4F6C}
\definecolor{matd48verde}{HTML}{2E8B57}

\titleformat{\section}{\normalfont\Large\bfseries\color{matd48azul}}{\thesection}{1em}{}
\titleformat{\subsection}{\normalfont\large\bfseries\color{matd48azul}}{\thesubsection}{1em}{}

\providecommand{\ListaTitulo}{Lista de Exercícios}

\setlength{\headheight}{14pt}
\pagestyle{fancy}
\fancyhf{}
\lhead{\small MATD48 --- Planejamento de Experimentos A}
\rhead{\small \ListaTitulo}
\cfoot{\thepage}
\renewcommand{\headrulewidth}{0.4pt}

\lstset{
  basicstyle=\ttfamily\small,
  breaklines=true,
  frame=single,
  columns=fullflexible,
  backgroundcolor=\color{gray!8},
  commentstyle=\color{matd48verde},
  keywordstyle=\color{matd48azul}\bfseries,
  language=R,
  extendedchars=true,
  literate=%
    {á}{{\'a}}1 {à}{{\`a}}1 {â}{{\^a}}1 {ã}{{\~a}}1
    {é}{{\'e}}1 {ê}{{\^e}}1 {í}{{\'i}}1 {ó}{{\'o}}1
    {ô}{{\^o}}1 {õ}{{\~o}}1 {ú}{{\'u}}1 {ç}{{\c c}}1
    {Á}{{\'A}}1 {À}{{\`A}}1 {Â}{{\^A}}1 {Ã}{{\~A}}1
    {É}{{\'E}}1 {Ê}{{\^E}}1 {Í}{{\'I}}1 {Ó}{{\'O}}1
    {Ô}{{\^O}}1 {Õ}{{\~O}}1 {Ú}{{\'U}}1 {Ç}{{\c C}}1
}

\hypersetup{colorlinks=true, linkcolor=matd48azul, urlcolor=matd48azul, citecolor=matd48azul}

\newcommand{\cabecalho}[2]{%
  \begin{center}
    {\Large\bfseries\color{matd48azul} #1}\\[0.3em]
    {\large #2}\\[0.3em]
    \rule{\linewidth}{0.6pt}
  \end{center}
  \vspace{0.5em}
}

\newenvironment{questao}[1]{\vspace{1em}\noindent\textbf{Questão #1.}\hspace{0.4em}}{\par}
\newenvironment{solucao}{\vspace{0.3em}\noindent\textit{Solução.}\hspace{0.4em}\color{black!85}}{\par\vspace{0.5em}}


\usepackage[utf8]{inputenc}
\usepackage[T1]{fontenc}
\usepackage[brazil]{babel}
\usepackage{fullpage}
\usepackage{amsmath,amssymb}
\usepackage{enumitem}
\usepackage{xcolor}
\usepackage{fancyhdr}
\usepackage{titlesec}
\usepackage{listings}
\usepackage{hyperref}
\usepackage{booktabs}
\usepackage{graphicx}

\definecolor{matd48azul}{HTML}{0B4F6C}
\definecolor{matd48verde}{HTML}{2E8B57}

\titleformat{\section}{\normalfont\Large\bfseries\color{matd48azul}}{\thesection}{1em}{}
\titleformat{\subsection}{\normalfont\large\bfseries\color{matd48azul}}{\thesubsection}{1em}{}

\providecommand{\ListaTitulo}{Lista de Exercícios}

\setlength{\headheight}{14pt}
\pagestyle{fancy}
\fancyhf{}
\lhead{\small MATD48 --- Planejamento de Experimentos A}
\rhead{\small \ListaTitulo}
\cfoot{\thepage}
\renewcommand{\headrulewidth}{0.4pt}

\lstset{
  basicstyle=\ttfamily\small,
  breaklines=true,
  frame=single,
  columns=fullflexible,
  backgroundcolor=\color{gray!8},
  commentstyle=\color{matd48verde},
  keywordstyle=\color{matd48azul}\bfseries,
  language=R,
  extendedchars=true,
  literate=%
    {á}{{\'a}}1 {à}{{\`a}}1 {â}{{\^a}}1 {ã}{{\~a}}1
    {é}{{\'e}}1 {ê}{{\^e}}1 {í}{{\'i}}1 {ó}{{\'o}}1
    {ô}{{\^o}}1 {õ}{{\~o}}1 {ú}{{\'u}}1 {ç}{{\c c}}1
    {Á}{{\'A}}1 {À}{{\`A}}1 {Â}{{\^A}}1 {Ã}{{\~A}}1
    {É}{{\'E}}1 {Ê}{{\^E}}1 {Í}{{\'I}}1 {Ó}{{\'O}}1
    {Ô}{{\^O}}1 {Õ}{{\~O}}1 {Ú}{{\'U}}1 {Ç}{{\c C}}1
}

\hypersetup{colorlinks=true, linkcolor=matd48azul, urlcolor=matd48azul, citecolor=matd48azul}

\newcommand{\cabecalho}[2]{%
  \begin{center}
    {\Large\bfseries\color{matd48azul} #1}\\[0.3em]
    {\large #2}\\[0.3em]
    \rule{\linewidth}{0.6pt}
  \end{center}
  \vspace{0.5em}
}

\newenvironment{questao}[1]{\vspace{1em}\noindent\textbf{Questão #1.}\hspace{0.4em}}{\par}
\newenvironment{solucao}{\vspace{0.3em}\noindent\textit{Solução.}\hspace{0.4em}\color{black!85}}{\par\vspace{0.5em}}


\begin{document}

\cabecalho{Gabarito --- Lista de Exercícios 04}{Delineamento completamente aleatorizado: modelo, ANOVA, intervalos de confiança e tamanho amostral (Capítulo 3 / Aula 04)}

\begin{questao}{1} \textbf{(Psicologia --- ANOVA a partir de somas de quadrados)}
\end{questao}
\begin{solucao}
\begin{enumerate}[label=(\alph*)]
  \item Com $t=3$, $N=30$: $\mathrm{g.l._{Trat}} = 2$, $\mathrm{g.l._{E}} = 27$,
    $\mathrm{g.l._{Total}} = 29$. $SQ_E = 402{,}6 - 148{,}2 = 254{,}4$.
    $MS_{\text{Trat}} = 148{,}2/2 = 74{,}1$; $MS_E = 254{,}4/27 \approx 9{,}42$.
    $F = 74{,}1/9{,}42 \approx \textbf{7{,}86}$.

  \begin{center}
  \begin{tabular}{lccccc}
  \toprule
  Fonte & g.l. & SQ & QM & $F$ \\
  \midrule
  Técnica & 2 & 148{,}2 & 74{,}10 & 7{,}86 \\
  Erro & 27 & 254{,}4 & 9{,}42 & \\
  Total & 29 & 402{,}6 & & \\
  \bottomrule
  \end{tabular}
  \end{center}

  \item Sim, diferem: $F=7{,}86 > F_{2,27;0{,}05}\approx 3{,}35$, logo rejeitamos
    $H_0:\mu_1=\mu_2=\mu_3$ a $\alpha=0{,}05$. Há evidência de que ao menos uma técnica difere das
    demais quanto à redução média de cortisol.
  \item A ANOVA global só permite afirmar que \emph{existe} alguma diferença entre as três médias
    --- \textbf{não} identifica quais pares diferem, nem a direção ou magnitude de cada diferença
    específica. Para isso são necessários procedimentos de comparação múltipla (contrastes,
    Tukey, etc.), fora do escopo desta lista.
\end{enumerate}
\end{solucao}

\begin{questao}{2} \textbf{(Agricultura --- intervalos de confiança)}
\end{questao}
\begin{solucao}
Com $MS_E=18{,}4$, $df_E=28$, $n=8$ por dose, $t_{28;0{,}025}\approx 2{,}048$.
\begin{enumerate}[label=(\alph*)]
  \item $\bar Y_{80} \pm t_{28;0{,}025}\sqrt{MS_E/8} = 74{,}3 \pm 2{,}048\sqrt{2{,}3} = 74{,}3 \pm
    3{,}11$, ou seja, $\textbf{IC}_{95\%} \approx (71{,}19;\ 77{,}41)$ sc/ha.
  \item $(\bar Y_{80}-\bar Y_{40}) \pm t_{28;0{,}025}\sqrt{MS_E(1/8+1/8)} = 5{,}6 \pm
    2{,}048\sqrt{4{,}6} = 5{,}6 \pm 4{,}39$, ou seja, $\textbf{IC}_{95\%}\approx (1{,}21;\ 9{,}99)$
    sc/ha --- como o intervalo não contém zero, há evidência de que a dose de 80 kg/ha produz
    maior produtividade média que a de 40 kg/ha.
  \item Porque a ANOVA assume variância comum $\sigma^2$ entre todos os $t$ tratamentos
    (homocedasticidade); combinar (\emph{pool}) a variabilidade das quatro doses em uma única
    estimativa $MS_E$, com $28$ graus de liberdade em vez de apenas $7$ por dose, produz uma
    estimativa de $\sigma^2$ mais precisa e, portanto, intervalos mais estreitos --- desde que a
    suposição de variância comum seja razoável (Aula 06).
\end{enumerate}
\end{solucao}

\begin{questao}{3} \textbf{(Ciência de dados --- notação matricial do DCA)}
\end{questao}
\begin{solucao}
\begin{enumerate}[label=(\alph*)]
  \item $\mathbf{X}$ é $12\times 3$. Como as três colunas de $\mathbf{X}$ são indicadoras de
    grupos disjuntos, $\mathbf{X}'\mathbf{X} = \mathrm{diag}(n_1,n_2,n_3) =
    \mathrm{diag}(4,3,5)$ (uma matriz $3\times 3$), e portanto
    $(\mathbf{X}'\mathbf{X})^{-1} = \mathrm{diag}(1/4,\,1/3,\,1/5)$.
  \item Porque $\mathbf{X}$ tem posto coluna completo ($\mathrm{posto}(\mathbf{X})=3=$ número de
    colunas): não há combinação linear não trivial das colunas de $\mathbf{X}$ que seja o vetor
    nulo, logo $\mathbf{X}'\mathbf{X}$ é invertível e \emph{toda} função linear de $\boldsymbol\beta$
    --- em particular cada $\mu_i$ isoladamente --- é estimável de forma única. Isso contrasta com
    a parametrização sobreparametrizada $\mu+\tau_i$ (posto $t-1$ para os $\tau_i$ dado $\mu$), em
    que só contrastes (como $\mu_i-\mu_k$) são estimáveis, não $\mu$ e $\tau_i$ isoladamente.
  \item $\mathrm{Var}(\hat\mu_2-\hat\mu_3) = \mathrm{Var}(\hat\mu_2)+\mathrm{Var}(\hat\mu_3)
    -2\,\mathrm{Cov}(\hat\mu_2,\hat\mu_3) = \sigma^2/n_2 + \sigma^2/n_3 - 0 =
    \sigma^2\!\left(\dfrac{1}{3}+\dfrac{1}{5}\right) = \dfrac{8}{15}\sigma^2$, já que
    $(\mathbf{X}'\mathbf{X})^{-1}$ diagonal implica $\mathrm{Cov}(\hat\mu_2,\hat\mu_3)=0$.
\end{enumerate}
\end{solucao}

\begin{questao}{4} \textbf{(Dedução --- $\hat{\boldsymbol\beta}$ é BLUE no DCA)}
\end{questao}
\begin{solucao}
\begin{enumerate}[label=(\alph*)]
  \item $E[\hat{\boldsymbol\beta}] = E\big[(\mathbf{X}'\mathbf{X})^{-1}\mathbf{X}'\mathbf{Y}\big]
    = (\mathbf{X}'\mathbf{X})^{-1}\mathbf{X}'E[\mathbf{Y}]
    = (\mathbf{X}'\mathbf{X})^{-1}\mathbf{X}'\mathbf{X}\boldsymbol\beta = \boldsymbol\beta$, usando
    $E[\mathbf{Y}]=\mathbf{X}\boldsymbol\beta$ (pois $E[\boldsymbol\varepsilon]=\mathbf 0$) e que
    $(\mathbf{X}'\mathbf{X})^{-1}\mathbf{X}'\mathbf{X}=\mathbf{I}_t$. Logo $\hat{\boldsymbol\beta}$
    é não viesado.
  \item $\mathrm{Cov}(\hat{\boldsymbol\beta}) =
    (\mathbf{X}'\mathbf{X})^{-1}\mathbf{X}'\,\mathrm{Cov}(\mathbf{Y})\,\mathbf{X}(\mathbf{X}'\mathbf{X})^{-1}
    = (\mathbf{X}'\mathbf{X})^{-1}\mathbf{X}'(\sigma^2\mathbf{I}_N)\mathbf{X}(\mathbf{X}'\mathbf{X})^{-1}
    = \sigma^2(\mathbf{X}'\mathbf{X})^{-1}\underbrace{\mathbf{X}'\mathbf{X}(\mathbf{X}'\mathbf{X})^{-1}}_{=\mathbf I_t}
    = \sigma^2(\mathbf{X}'\mathbf{X})^{-1}$. Como $\mathbf{X}'\mathbf{X}=\mathrm{diag}(n_1,\ldots,n_t)$
    para a matriz de indicadoras (colunas ortogonais entre si), sua inversa é
    $\mathrm{diag}(1/n_1,\ldots,1/n_t)$, logo $\mathrm{Cov}(\hat{\boldsymbol\beta}) =
    \mathrm{diag}(\sigma^2/n_1,\ldots,\sigma^2/n_t)$: diagonal, $i$-ésimo elemento $\sigma^2/n_i$.
  \item \textbf{Teorema (Gauss-Markov).} Sob $E[\boldsymbol\varepsilon]=\mathbf 0$ e
    $\mathrm{Cov}(\boldsymbol\varepsilon)=\sigma^2\mathbf{I}_N$ (sem exigir normalidade), se
    $\boldsymbol\lambda'\boldsymbol\beta$ é estimável, então $\boldsymbol\lambda'\hat{\boldsymbol\beta}$
    é o BLUE de $\boldsymbol\lambda'\boldsymbol\beta$: entre todos os estimadores lineares
    $\mathbf{c}'\mathbf{Y}$ não viesados para $\boldsymbol\lambda'\boldsymbol\beta$, ele tem a
    menor variância. As hipóteses do teorema --- média zero e covariância $\sigma^2\mathbf I_N$ dos
    erros --- são exatamente as assumidas no enunciado, e $\mathbf{X}$ tem posto completo, o que
    torna \emph{todo} $\boldsymbol\lambda'\boldsymbol\beta$ estimável (item b). Em particular,
    tomando $\boldsymbol\lambda = \mathbf{e}_i$ (o $i$-ésimo vetor canônico), $\mu_i =
    \mathbf{e}_i'\boldsymbol\beta$ é estimável e $\hat\mu_i = \mathbf{e}_i'\hat{\boldsymbol\beta} =
    \bar Y_{i\cdot}$ é o BLUE de $\mu_i$: nenhum outro estimador linear não viesado de $\mu_i$ tem
    variância menor que $\sigma^2/n_i$.
\end{enumerate}
\end{solucao}

\begin{questao}{5} \textbf{(Número de réplicas --- ciência de dados)}
\end{questao}
\begin{solucao}
\begin{enumerate}[label=(\alph*)]
  \item O menor $n$ tabelado com poder $\geq 0{,}80$ é $\textbf{n=40}$ por grupo (poder $0{,}80$
    exatamente), totalizando $3\times 40 = \textbf{120}$ observações.
  \item Com $n=20$, o poder tabelado é $0{,}50$: se a diferença real entre as estratégias de
    \emph{cache} for de fato do tamanho assumido no piloto ($f\approx 0{,}30$), a equipe teria
    apenas cerca de 50\% de chance de detectá-la como estatisticamente significativa --- ou seja,
    quase a mesma chance de um "cara ou coroa" de o experimento terminar inconclusivo mesmo
    havendo efeito real (erro Tipo II de $\beta \approx 0{,}50$).
  \item \textbf{Subestimar} $\sigma$ tornaria o planejamento otimista demais: um $\sigma$ menor do
    que o real infla artificialmente o $f$ de Cohen calculado a partir das médias piloto (pois
    $f=\sigma_\mu/\sigma$), fazendo a tabela superestimar o poder alcançado para um dado $n$ --- na
    prática, o poder real seria menor que o nominal, e o experimento ficaria subdimensionado. Por
    isso pilotos muito pequenos, que tendem a estimar $\sigma$ de forma instável, pedem cautela
    (ou uma margem de segurança) no planejamento amostral.
\end{enumerate}
\end{solucao}

\end{document}
